\documentclass[11pt,a4paper]{article}
\usepackage[T1]{fontenc}
\usepackage[utf8]{inputenc}
\usepackage{lmodern}
\usepackage{geometry}
\geometry{left=27mm,right=27mm,top=25mm,bottom=26mm}
\usepackage{amsmath,amssymb,amsthm,mathtools,bm}
\usepackage{graphicx,booktabs,array,microtype}
\usepackage{hyperref}
\hypersetup{colorlinks=true,linkcolor=blue!45!black,citecolor=blue!45!black,urlcolor=blue!55!black}
\usepackage{xcolor}
\usepackage{enumitem}
\setlist{nosep,leftmargin=*}
\usepackage{caption}
\usepackage{authblk}
\renewcommand{\Authfont}{\normalsize}
\renewcommand{\Affilfont}{\small}
\setlength{\affilsep}{0.5em}
\captionsetup{font=small}
\newtheorem{theorem}{Theorem}[section]
\newtheorem{proposition}[theorem]{Proposition}
\newtheorem{lemma}[theorem]{Lemma}
\newtheorem{corollary}[theorem]{Corollary}
\theoremstyle{definition}\newtheorem{definition}[theorem]{Definition}
\theoremstyle{remark}\newtheorem{remark}[theorem]{Remark}
\DeclareMathOperator{\diag}{diag}
\DeclareMathOperator{\cond}{cond}
\DeclareMathOperator{\spec}{spec}
\DeclareMathOperator{\tr}{tr}
\newcommand{\C}{\mathbb{C}}
\newcommand{\Rn}{\mathbb{R}}
\newcommand{\norm}[1]{\left\lVert #1\right\rVert}
\newcommand{\abs}[1]{\left\lvert #1\right\rvert}
\newcommand{\ee}{\mathrm{e}}
\newcommand{\op}{\mathrm{op}}
\newcommand{\I}{\mathbb{I}}
\newcommand{\ket}[1]{\lvert #1\rangle}
\newcommand{\sigmin}{\sigma_{\min}}
\newcommand{\kg}{\kappa_2}
\title{\bfseries Metric Conditioning and Stable Spectral Inversion\\[3pt]
 near Exceptional Points in Pseudo-Hermitian Systems}
\author[1]{Dorcas Attuabea Addo}
\author[2]{Richard Kena Boadi}
\author[2]{Jeffrey Ezrean}
\author[3]{Richard Kwame Ansah}
\author[4]{Peter Yeboah}
\affil[1]{Department of Mathematics Education, University of Education, Winneba, Ghana}
\affil[2]{Department of Mathematics, Kwame Nkrumah University of Science and Technology, Kumasi, Ghana}
\affil[3]{Department of Mathematics and Statistics, University of Energy and Natural Resources, Sunyani, Ghana}
\affil[4]{Department of Teacher Education, Kwame Nkrumah University of Science and Technology, Kumasi, Ghana}
\date{}
\begin{document}
\maketitle
\begin{abstract}
For a diagonalizable non-Hermitian Hamiltonian with real spectrum, a positive metric operator
can render the Hamiltonian self-adjoint in a modified inner product.  Near an exceptional
point, however, the metric may become strongly ill-conditioned even when the Hamiltonian
remains uniformly invertible in the reference Euclidean norm.  We quantify this separation.
After deriving general metric-based perturbation bounds, we analyze two exactly solvable
families.  In a two-level model the condition-number-minimizing metric yields a Euclidean
invertibility certificate that is sharp.  In contrast, for a coupled three-level family
$A_\varepsilon$ approaching a defective but nonzero eigenvalue, the Euclidean inverse remains
uniformly bounded while every compatible positive metric becomes ill-conditioned.  We prove
the sharp asymptotic law
\[
 \kappa_*(A_\varepsilon)\sim \frac{8}{\varepsilon^2},\qquad \varepsilon\downarrow0,
\]
where $\kappa_*$ is the minimum spectral condition number over all positive compatible
metrics, and we construct an asymptotically optimal metric attaining the constant $8$.
Consequently, the best metric-induced Euclidean invertibility certificate vanishes linearly,
although the exact distance to singularity stays bounded away from zero.  A normalized-solution
perturbation estimate shows that inversion itself remains stable.  The same separation is then
embedded in a genuinely coupled $N$-level bidiagonal family, uniformly in dimension.  A standard
Hermitian dilation is included only to emphasize that singular-value conditioning need not
inherit the metric singularity.  These results distinguish exceptional-point eigenvector
geometry from ordinary inversion stability and identify a precise limitation of metric-based
conditioning as a proxy for Euclidean spectral inversion.
\end{abstract}
\noindent\textbf{Keywords:} pseudo-Hermitian Hamiltonian; quasi-Hermiticity; metric operator;
exceptional point; spectral inversion; pseudospectrum; eigenvector coalescence.\\
\noindent\textbf{MSC 2020 (provisional):} 47A10, 81Q10, 65F35.
